% GENERATED FILE. Edit canonical lessons, then run npm run generate:books.
% canonical-source-hash: fnv1a64:a648f7d11c5650d5
% canonical-lessons: UR-C21-kyun, UR-C21-kaun, UR-C21-kyunki, UR-C21-ki, UR-C21-ke-liye, UR-R21-why-and-what-for

\chapter{Why, and What For}
\label{ch:ur-why-and-what-for}
\hlchaptermodality{21}

\begin{chapteropening}
\textbf{By the end of this chapter:} \emph{I can ask why and who, answer why with a cause using \ur{کیونکہ} or with a purpose using \ur{کے} \ur{لیے}, and report what I think with \ur{کہ}.}

Urdu's k- question family is completed -- \ur{کیا}, \ur{کہاں}, \ur{کیوں}, \ur{کون} -- and the reader answers \ur{کیوں} two different ways, with a cause and with a purpose, then reports an opinion over \ur{سوچنا} from chapter seven.
\end{chapteropening}

\section[kyūṅ?]{\ur{کیوں} (kyūṅ) — why}

\label{lesson:UR-C21-kyun}



\begin{quote}
Chapter seventeen gathered Urdu's \textbf{k-} question words and named them a family. One member was left out, and it is the one that asks for a reason.

\begin{itemize}
\raggedright
  \item \emph{Recall:} the joiner two lessons back, the one that turns — \emph{lekin}
\end{itemize}
\end{quote}

\subsection*{You'll want to know: \ur{کیوں}}
\begin{quote}
\textbf{\ur{کیوں؟}} — \emph{kyūṅ?} — \textbf{why?}
\end{quote}

\begin{quote}
\textbf{\ur{آپ} \ur{کیوں} \ur{نہیں} …\ur{؟}} — \emph{āp kyūṅ nahīṅ āye?} — \textbf{why did you not come?}
\end{quote}

Every letter is one you own: \textbf{\ur{ک}}\ur{،} \textbf{\ur{ی}}\ur{،} \textbf{\ur{و}}\ur{،} and the humming \textbf{\ur{ں}} of chapter six. It stands where \textbf{\ur{کہاں}} stands — in front of the verb, not at the front of the sentence, which is where English puts \emph{why}.

\begin{cousinweb}
\noindent
\begin{tabularx}{\linewidth}{@{}>{\raggedright\arraybackslash}X>{\raggedright\arraybackslash}X@{}}
\toprule
\textbf{word} & \textbf{asks} \\
\midrule
\textbf{\ur{کیا}} & what \\
\textbf{\ur{کہاں}} & where \\
\textbf{\ur{کیوں}} & why \\
\bottomrule
\end{tabularx}

One stem, three questions. It is Sanskrit \textbf{ka-}, the interrogative stem behind Latin \emph{quis} and \emph{quid}, and behind English \textbf{who}, \textbf{what}, \textbf{where} and \textbf{why} — the \emph{wh-} of English and the \emph{k-} of Urdu are the same ancient consonant, moved apart by one sound change.

\textbf{\ur{کیوں}} on its own, said flatly, is a complete question, and it is what a child says twenty times a day in either language.
\end{cousinweb}

\subsection*{Your turn}
Say these aloud:

\begin{itemize}
\raggedright
  \item \emph{kyūṅ?} — on its own
  \item the three k- words in a row — \emph{kyā, kahāṅ, kyūṅ}
  \item \emph{āp kyūṅ nahīṅ āye?} — why did you not come?
\end{itemize}

\begin{tcolorbox}[breakable,colback=teal!4,colframe=teal!35!black,title={Before you move on}]
Which stem do \textbf{\ur{کیا}}\ur{،} \textbf{\ur{کہاں}} and \textbf{\ur{کیوں}} share? (\textbf{Sanskrit ka-}, the interrogative stem.) Which English consonant answers Urdu's \textbf{k-}? (\textbf{The \emph{wh-}} of \emph{who}, \emph{what}, \emph{where}, \emph{why}.) And where does \textbf{\ur{کیوں}} stand in the sentence? (\textbf{In front of the verb}, not at the front.)

Source: \href{https://www.unicode.org/charts/PDF/U0600.pdf}{Unicode Arabic chart}.
\end{tcolorbox}
\section[kaun?]{\ur{کون} (kaun) — who}

\label{lesson:UR-C21-kaun}



\begin{quote}
The chapter-eight writing lesson used a three-letter word to practise the letter \textbf{\ur{و}} and glossed it in passing. It never became a word you own.

\begin{itemize}
\raggedright
  \item \emph{Recall:} the question word from the lesson before — \emph{kyūṅ}
\end{itemize}
\end{quote}

\subsection*{You'll want to know: \ur{کون}}
\begin{quote}
\textbf{\ur{یہ} \ur{کون} \ur{ہے؟}} — \emph{yih kaun hai?} — \textbf{who is this?}
\end{quote}

\begin{quote}
\textbf{\ur{کون؟}} — \emph{kaun?} — \textbf{who?}
\end{quote}

Three letters, all long since taught: \textbf{\ur{ک}}\ur{،} \textbf{\ur{و}}\ur{،} \textbf{\ur{ن}}. The \textbf{\ur{و}} is the same letter chapter eight taught as a vowel, and this word is where you first met it without being told what the word meant.

\begin{cousinweb}
\textbf{\ur{کون}} is the fourth of the \textbf{k-} family, from Sanskrit \emph{kaḥ punaḥ}, \textquotedblleft{}who then,\textquotedblright{} worn to a syllable. It is the direct cousin of English \textbf{who} — the ancient word began with a \emph{k}-sound, and English wrote it \emph{hw-} and then \emph{wh-} while keeping the sound, which is why \emph{who} still starts with the breath the letter \emph{w} does not say.

Now the family is whole: \textbf{\ur{کیا}} what, \textbf{\ur{کہاں}} where, \textbf{\ur{کیوں}} why, \textbf{\ur{کون}} who. Four questions and one stem, and a reader who hears the \textbf{k-} knows a question is coming before the sentence ends.
\end{cousinweb}

\subsection*{Your turn}
Say these aloud:

\begin{itemize}
\raggedright
  \item \emph{kaun?} — on its own
  \item \emph{yih kaun hai?}
  \item and answer it — \emph{yih merī māṅ hai}
  \item the whole family — \emph{kyā, kahāṅ, kyūṅ, kaun}
\end{itemize}

\begin{tcolorbox}[breakable,colback=teal!4,colframe=teal!35!black,title={Before you move on}]
How many members does the \textbf{k-} family now have in this book, and what are they? (\textbf{Four} — \emph{kyā, kahāṅ, kyūṅ, kaun}.) Which English word is \textbf{\ur{کون}}'s direct cousin? (\textbf{Who}.) And which letter in the middle of it did chapter eight teach? (\textbf{\ur{و}}.)

Source: \href{https://www.unicode.org/charts/PDF/U0600.pdf}{Unicode Arabic chart}.
\end{tcolorbox}
\section[kyūṅki]{\ur{کیونکہ} (kyūṅki) — because}

\label{lesson:UR-C21-kyunki}



\begin{quote}
The question is now askable. The answer is not: nothing in this book can put a reason after a claim.

\begin{itemize}
\raggedright
  \item \emph{Recall:} the two items before this one — \emph{kyūṅ}, \emph{kaun}
\end{itemize}
\end{quote}

\subsection*{You'll want to know: \ur{کیونکہ}}
\begin{quote}
\textbf{\ur{پانی} \ur{نہیں} \ur{ہے} \ur{کیونکہ} …} — \emph{pānī nahīṅ hai kyūṅki …} — \textbf{there is no water, because …}
\end{quote}

And over the wellbeing frame of chapter four, whose adjective carries two letters still ahead of you in the script ladder: \emph{maiṅ ṭhīk nahīṅ hūṅ kyūṅki …}, \textbf{I am not well, because …}

Look at the word and the last lesson is inside it: \textbf{\ur{کیوں}} plus \textbf{\ur{کہ}}. Urdu builds \emph{because} by welding \emph{why} to \emph{that} — the language answers its own question in a single word.

Spanish does exactly this: \emph{por qué} asks, \emph{porque} answers, and the only difference is a space and an accent. English did it too, in \emph{forasmuch as}, before the phrase went out of use.

\begin{grammarlens}[title={Grammar Lens: claim first, reason after}]
\noindent
\begin{tabularx}{\linewidth}{@{}>{\raggedright\arraybackslash}X>{\raggedright\arraybackslash}X@{}}
\toprule
\textbf{head word} & \textbf{what follows it} \\
\midrule
\textbf{\ur{لیکن}} & an objection \\
\textbf{\ur{کیونکہ}} & a reason \\
\bottomrule
\end{tabularx}

Both stand at the head of the clause they govern, and both come \textbf{after} the claim they attach to. Put them in one sentence and you have the whole shape of an argument: \emph{ṭhīk hai, lekin … kyūṅki …}
\end{grammarlens}

\subsection*{Your turn}
Say these aloud:

\begin{itemize}
\raggedright
  \item the question, then the answer word — \emph{kyūṅ? … kyūṅki}
  \item \emph{maiṅ ṭhīk nahīṅ hūṅ kyūṅki maiṅ nahīṅ samajhtā}
  \item an objection and a reason in one breath — \emph{ṭhīk hai, lekin … kyūṅki …}
\end{itemize}

\begin{tcolorbox}[breakable,colback=teal!4,colframe=teal!35!black,title={Before you move on}]
What two words is \textbf{\ur{کیونکہ}} built from? (\textbf{\ur{کیوں}}, why, and \textbf{\ur{کہ}}, that.) Which Spanish pair is built the same way? (\emph{\textbf{Por qué}} and \emph{\textbf{porque}}.) And which comes first, the claim or the reason? (\textbf{The claim} — the reason follows it.)

Source: \href{https://www.unicode.org/charts/PDF/U0600.pdf}{Unicode Arabic chart}.
\end{tcolorbox}
\section[ki]{\ur{کہ} (ki) — that}

\label{lesson:UR-C21-ki}



\begin{quote}
\emph{sochnā}, \textquotedblleft{}to think,\textquotedblright{} has been in this book since chapter seven, and in all that time nobody has been able to say what they think.

\begin{itemize}
\raggedright
  \item \emph{Recall:} the reason word from the lesson before, and the two words inside it — \emph{kyūṅki}
\end{itemize}
\end{quote}

\subsection*{You'll want to know: \ur{کہ}}
\begin{quote}
\textbf{… \ur{کہ} \ur{یہ} \ur{پانی} \ur{ہے}} — \emph{maiṅ sochtā hūṅ ki yih pānī hai} — \textbf{I think that this is water.}
\end{quote}

The verb in front of \textbf{\ur{کہ}} is \emph{sochtā}, whose \emph{che} and \emph{te} are still ahead of you in the script ladder, so the frame is printed in romanization and the clause \textbf{\ur{کہ}} introduces is printed whole.

\textbf{\ur{کہ}} does one job: it takes a whole sentence and makes it the object of a verb of thinking, knowing or saying. Everything before it is the frame; everything after it is the thought, in exactly the order it would have had standing alone.

Two letters, both yours — \textbf{\ur{ک}} and \textbf{\ur{ہ}} — and you have been reading them together since the last lesson, welded onto \textbf{\ur{کیوں}}.

\begin{grammarlens}[title={Grammar Lens: what English drops and Urdu keeps}]
English can leave \emph{that} out — \emph{I think this is right} is ordinary English. Urdu cannot: \textbf{\ur{کہ}} is obligatory, and a sentence built without it is not merely informal, it is wrong.

The word is Persian \emph{ki}, borrowed rather than inherited — which is why it sits unstressed and one-shaped wherever it lands. Compare \textbf{\ur{کہاں}}, \emph{where}, which begins with the same two letters and is a completely different word from a different source: one of the places Urdu's two ancestries look identical on the page.
\end{grammarlens}

\subsection*{Your turn}
Say these aloud:

\begin{itemize}
\raggedright
  \item \emph{maiṅ sochtā hūṅ ki yih ṭhīk hai}
  \item \emph{maiṅ samajhtā hūṅ ki …}
  \item an opinion, then an objection to it — \emph{maiṅ sochtā hūṅ ki ṭhīk hai, lekin …}
\end{itemize}

\begin{tcolorbox}[breakable,colback=teal!4,colframe=teal!35!black,title={Before you move on}]
What does \textbf{\ur{کہ}} attach to the verb in front of it? (\textbf{A whole sentence}, as its object.) Can it be left out, the way English leaves out \emph{that}? (\textbf{No} — Urdu requires it.) And where have you already read these two letters? (\textbf{Inside \ur{کیونکہ}}.)

Source: \href{https://www.unicode.org/charts/PDF/U0600.pdf}{Unicode Arabic chart}.
\end{tcolorbox}
\section[ke liye]{\ur{کے} \ur{لیے} (ke liye) — for, in order to}

\label{lesson:UR-C21-ke-liye}



\begin{quote}
Chapter six taught the \textbf{-\ur{نا}} ending every Urdu verb wears in the dictionary. Fourteen verbs have that ending and not one of them has ever been the reason for anything.

\begin{itemize}
\raggedright
  \item \emph{Recall:} the complementiser from the lesson before — \emph{ki}
\end{itemize}
\end{quote}

\subsection*{You'll want to know: \ur{کے} \ur{لیے}}
\begin{quote}
\textbf{\ur{پانی} \ur{کے} \ur{لیے}} — \emph{pānī ke liye} — \textbf{for water.}
\end{quote}

\begin{quote}
\emph{paṛhne ke liye} — \textbf{in order to read.}
\end{quote}

Two words, five letters, and every one of them long since yours: \textbf{\ur{ک}}\ur{،} \textbf{\ur{ے}}\ur{،} \textbf{\ur{ل}}\ur{،} \textbf{\ur{ی}}\ur{،} \textbf{\ur{ے}}. Put it after a noun and it means \textbf{for}. Put it after a verb and it means \textbf{in order to} — and the verb's \textbf{-\ur{نا}} becomes \textbf{-\ur{نے}} first, which is the one change the construction asks for.

\begin{grammarlens}[title={Grammar Lens: two ways to say why}]
\noindent
\begin{tabularx}{\linewidth}{@{}>{\raggedright\arraybackslash}X>{\raggedright\arraybackslash}X@{}}
\toprule
\textbf{shape} & \textbf{says} \\
\midrule
\textbf{\ur{کیونکہ}} + a whole clause & because something is so \\
verb + \textbf{\ur{کے} \ur{لیے}} & in order to do something \\
\bottomrule
\end{tabularx}

The first gives a cause, the second gives a purpose, and a learner who owns both can answer \emph{kyūṅ?} two different ways.

\textbf{\ur{لیے}} is the perfective participle of \textbf{\ur{لینا}}, \textquotedblleft{}to take,\textquotedblright{} which chapter eight taught. Urdu builds \emph{for} out of \emph{having taken}: a purpose is what the action was taken for. English went the other way and built \emph{for} out of a word for \emph{before}.
\end{grammarlens}

\subsection*{Your turn}
Say these aloud:

\begin{itemize}
\raggedright
  \item \emph{pānī ke liye}
  \item the infinitive, then the purpose — \emph{paṛhnā … paṛhne ke liye}
  \item \emph{samajhne ke liye} — in order to understand
  \item both answers to \emph{kyūṅ?} — \emph{kyūṅki …} and \emph{… ke liye}
\end{itemize}

\begin{tcolorbox}[breakable,colback=teal!4,colframe=teal!35!black,title={Before you move on}]
What does the infinitive's \textbf{-\ur{نا}} become before \textbf{\ur{کے} \ur{لیے}}? (\textbf{-\ur{نے}}.) Which verb is \textbf{\ur{لیے}} built on? (\textbf{\ur{لینا}}, to take — chapter eight.) And what does \textbf{\ur{کے} \ur{لیے}} give that \textbf{\ur{کیونکہ}} does not? (\textbf{A purpose}, rather than a cause.)

Source: \href{https://www.unicode.org/charts/PDF/U0600.pdf}{Unicode Arabic chart}.
\end{tcolorbox}
\section[kyūṅ · kyūṅki · ki]{Asking why, and answering it two ways}

\label{lesson:UR-R21-why-and-what-for}



\begin{quote}
Five items landed in this chapter. Two ask, two answer, and one hangs a sentence off a verb. Name all five.
\end{quote}

\subsection*{Your turn}
\begin{itemize}
\raggedright
  \item \emph{Say it:} the whole k- family — \emph{kyā, kahāṅ, kyūṅ, kaun}
  \item \emph{Say it:} an opinion — \emph{maiṅ sochtā hūṅ ki yih ṭhīk hai}
  \item \emph{Say it:} a cause, then a purpose — \emph{kyūṅki …} … \emph{samajhne ke liye}
  \item \emph{Recall:} six lessons back, the joiner that turns — \emph{lekin}
  \item \emph{Recall:} nine lessons back, the joiner that chooses — \emph{yā}
\end{itemize}

\begin{tcolorbox}[breakable,colback=teal!4,colframe=teal!35!black,title={Before you move on}]
Which of this chapter's words answers \emph{kyūṅ?} with a cause, and which with a purpose? (\textbf{\ur{کیونکہ}} with a cause; \textbf{\ur{کے} \ur{لیے}} with a purpose.) Which one is obligatory where English drops its equivalent? (\textbf{\ur{کہ}}.) And what is inside \textbf{\ur{کیونکہ}}? (\textbf{\ur{کیوں}} and \textbf{\ur{کہ}}, the chapter's own two other items.)
\end{tcolorbox}
